% Results section draft -- Project Sidewalk disagreement study.
% Every number here is produced by sidewalk/manifest.py and sidewalk/decompose.py
% over 12 city instances. Regenerate before submission; the platforms are live
% and the counts drift.

\section{Results}
\label{sec:results}

Our corpus is every label from the twelve live Project Sidewalk city
deployments: 811{,}368 labels in total, of which 177{,}599 carry at least three
independent reviews. Those reviewed labels rest on 746{,}975 individual review
decisions contributed by 11{,}480 reviewers.

\subsection{Two Kinds of Disagreement}

We decompose each label's review record into two quantities. \emph{Split}
measures how evenly the reviewers who took a position divided between agree and
disagree; \emph{unsure} measures the share who declined to take one. A label can
score high on either, both, or neither.

Both populations are substantial: reviewers split into opposing camps on
23{,}638 labels (13.3\%) and answered ``not sure'' at a rate of one third or
more on 7{,}299 (4.1\%). The two are only weakly related
($\rho = 0.14$), which is the first evidence that they are distinct phenomena
rather than two readings of a single underlying difficulty.

Disagreement is not distributed uniformly across the label taxonomy. Missing
curb ramps draw the highest rate of split reviews, consistent with a judgement
about whether a ramp \emph{ought} to exist rather than about what is visible.
Obstacles draw the highest rate of ``not sure'' responses. Split reviews also
rise monotonically with the severity rating assigned by the original labeller,
which is what a disagreement about where a threshold sits would predict.

\subsection{Where the Disagreement Lives}

Reviewers are not interchangeable. Among the 3{,}333 reviewers with at least
twenty reviews, ``not sure'' rates range from 0\% to 70.0\% with a median of
3.0\%. A raw count of unsure votes on a label therefore confounds a property of
the sidewalk with a property of whoever happened to review it.

We separate the two by fitting an additive model
$y_{ij} \sim \mu + a_i + b_j$ over individual review decisions, where $a_i$ is a
label effect and $b_j$ a reviewer effect, estimated by alternating means and
shrunk toward zero in proportion to the number of reviews supporting each label.

The model must be fitted \emph{within} each city. Deployments differ in both
their reviewer populations and their base rates, so a pooled fit absorbs
between-city variation into the reviewer term and materially overstates it: the
pooled estimate attributes 20.8\% of the variance in ``not sure'' responses to
reviewers against 14.3\% to labels, whereas the within-city estimates reverse
that ordering in ten of twelve deployments.

Fitted correctly, both outcomes are label-driven, and they differ in degree
rather than in kind (Table~\ref{tab:variance}). Disagreement is overwhelmingly
a property of the label, exceeding the reviewer term by a factor of roughly
five. Uncertainty is only weakly so: reviewer identity accounts for about twice
as much of the variance in ``not sure'' as it does in disagreement, and the
label for about half as much.

\begin{table}[tb]\centering
\caption{Variance in individual review decisions attributable to the label
versus the reviewer, estimated within each city (median across twelve
deployments). Disagreement is label-dominated in 12/12 cities; uncertainty in
10/12.}
\label{tab:variance}
\begin{tabular}{lcc}
\hline
\textbf{Source} & \textbf{Disagree} & \textbf{Unsure} \\
\hline
Label     & \textbf{52.0\%} & 27.1\% \\
Reviewer  & 9.7\%  & 20.7\% \\
label:reviewer & \textbf{5.3:1} & 1.3:1 \\
\hline
\end{tabular}
\end{table}

The practical consequence is that raw and reviewer-adjusted estimates of which
labels are ``unclear'' correlate at only $\rho = 0.68$: removing reviewer
tendency substantially reorders the population.

\subsection{More Review Does Not Produce Consensus}

Because the review records preserve each reviewer's identity and position in the
sequence, we can compare a label's first three reviews against its later ones
\emph{within the same label}, which removes any selection effect arising from
which labels attract additional review. We restrict to labels with at least six
reviews and subtract each reviewer's estimated tendency, fitted within their own
city, before comparing.

The result is a null, and it is the point. Neither form of disagreement declines
over a label's lifetime. Uncertainty rises slightly and consistently: the
adjusted delta is positive in eight of eleven deployments with sufficient data,
with a median of $+0.008$. Disagreement shows no consistent direction at all,
positive in three of eleven and with a median of $-0.005$
(Table~\ref{tab:resolve}).

Two analyses that we report but do not rely on illustrate why this had to be
done carefully. Without adjusting for reviewer tendency, uncertainty appears to
\emph{decline} over a label's lifetime ($0.106 \rightarrow 0.096$); that
apparent resolution is entirely explained by early reviewers being more
unsure-prone than later ones (mean reviewer effect $+0.043$ versus $+0.020$).
And when reviewer effects are estimated by pooling across deployments rather
than within them, disagreement appears to rise sharply ($+0.026$); that
disappears once each city's reviewers are modelled separately. Pooled estimates
in this corpus absorb between-city differences into the reviewer term, and both
of these artefacts point the same way---toward a tidier story than the data
supports.

\begin{table}[tb]\centering
\caption{Within-label comparison of a label's first three reviews against its
later ones, reviewer-adjusted within city. Delta is later minus early; positive
means the label became \emph{more} contested. Neither outcome declines.}
\label{tab:resolve}
\begin{tabular}{lccc}
\hline
\textbf{Outcome} & \textbf{Median $\Delta$} & \textbf{Positive in} & \textbf{Cities} \\
\hline
Unsure   & $+0.008$ & 8/11 & consistent \\
Disagree & $-0.005$ & 3/11 & no direction \\
\hline
\end{tabular}
\end{table}

The practical reading is narrow but sufficient for our argument: additional
review does not resolve these labels. Whatever the disagreement is, it is not
the kind of thing that more eyes settle.

\subsection{Quality Control Removes the Contested Cases}

Project Sidewalk's published quality criterion retains a label when it has more
than two agreements and no more than two disagreements. Applied to our corpus it
retains 54.5\% of reviewed labels overall. Applied to the contested populations
it retains 8.6\% of split labels and 6.8\% of unsure ones---discarding 91\% and
93\% respectively.

The filter is therefore not removing noise uniformly. It is removing, almost
completely, the subset of the record in which human judgement about
accessibility diverged. What survives into the analysable dataset is the portion
on which agreement was never in question.

This is not an artefact of any one deployment. Contested labels are retained at
a lower rate than the city average in all twelve cities, and leaving any single
city out moves the pooled figures by at most 1.3 percentage points---the largest
shift, to 9.9\% and 7.5\%, comes from excluding Seattle, which alone
contributes 30\% of the reviewed corpus. Bootstrap intervals over labels are
narrow: 8.6\% [8.3, 9.0] for split labels and 6.8\% [6.2, 7.3] for unsure ones,
against 54.5\% [54.2, 54.7] overall.

% ---------------------------------------------------------------------------
% TODO before submission
%   - regenerate all numbers; the city deployments are live
%   - DONE per-city breakdown (sidewalk/by_city.py); QC finding holds 12/12
%   - DONE bootstrap CIs on the QC retention rates
%   - DONE resolution test per city; the pooled 'disagreement worsens' result
%     did not replicate (3/11 cities) and has been removed from the claim
%   - verify that vote_index is chronological with the maintainers; the
%     within-label test assumes it, and the only current support is that
%     prolific reviewers appear at lower indices (rho = -0.16)
%   - state the additive-model assumption explicitly; a crossed random-effects
%     logistic fit is the stricter version if a reviewer objects
% ---------------------------------------------------------------------------
